\chapter{Implementation and Experimental Setup}
\label{chap:implementation}

This chapter describes the implementation of the experimental pipeline and the
evaluation framework. The implementation is intentionally script-oriented:
data preparation, model training, explanation, fairness measurement, and
bootstrap analysis are separate stages that share a single configuration.

\section{Technology Stack}

\begin{table}[H]
\centering
\caption{Technology stack used in the experiments.}
\label{tab:technology-stack}
\begin{tabular}{lll}
\toprule
\textbf{Component} & \textbf{Technology} & \textbf{Purpose}\\
\midrule
Language & Python & Experiment orchestration\\
Data processing & pandas, NumPy & Cleaning and tabular data\\
Model & LightGBM & Gradient-boosted trees\\
Explainability & SHAP / TreeSHAP & Global and local attribution\\
Metrics & scikit-learn & Predictive and ranking metrics\\
Statistics & SciPy & Correlation and bootstrap support\\
Document & LaTeX & Reproducible thesis output\\
\bottomrule
\end{tabular}
\end{table}

Python was selected because the required machine-learning and statistical
libraries provide compatible implementations of the complete pipeline.
LightGBM was preferred to a generic black-box model because it supports
categorical and missing values and exposes the built-in importance measures
required by RQ1. LaTeX was used to keep tables, equations, cross-references,
and bibliography generation reproducible.

\section{System Implementation}

The implementation is organized as a deterministic pipeline. The data-loading
stage validates the schema and removes duplicates and logically impossible
records. The preparation stage creates the age-group and joint
group-by-outcome stratification labels, encodes categorical variables, and
returns the training and test partitions. The training stage fits M0 and M1
with a shared parameter object. A separate training path supplies reweighing
weights for M0-RW and M1-RW.

The evaluation stage receives fitted models and the same test partition. It
calculates predictive metrics, age-group metrics, built-in importance, SHAP
values, and bootstrap intervals. Keeping evaluation outside the training
function prevents a model variant from silently using a different threshold or
test set. Every comparison records the model condition, split seed, feature
set, and whether sample weights were used.

\section{Implementation Challenges}

The first challenge was distinguishing invalid observations from legitimate
distributional extremes. A generic z-score filter removed every applicant
aged 47 or older, so it was replaced with explicit logical checks for impossible
age and employment-length values. The second challenge was avoiding leakage:
missing-value imputation and feature transformations are not fitted before the
split, and the test set is used only for evaluation. The third challenge was
preserving the meaning of categorical variables. Arbitrary integer ordering was
replaced with LightGBM's native categorical handling.

\section{Test Set Construction}

The primary test set contains 6,482 applicants: 2,444 under 25 and 4,038 aged
25 and above. It is not a manually authored benchmark; it is the held-out
portion of the cleaned public dataset. Stratification on age group and outcome
preserves all four group-by-outcome cells. The test set is never used for
hyperparameter selection or preprocessing. Ten additional stratified splits
are used as a robustness check.

\begin{table}[H]
\centering
\caption{Primary test-set composition.}
\label{tab:test-composition}
\begin{tabular}{lrrr}
\toprule
\textbf{Age group} & \textbf{All} & \textbf{Good} & \textbf{Bad}\\
\midrule
Under 25 & 2,444 & 1,874 & 570\\
25 and above & 4,038 & 3,191 & 847\\
\textbf{Total} & \textbf{6,482} & \textbf{5,065} & \textbf{1,417}\\
\bottomrule
\end{tabular}
\end{table}

\section{Evaluation Metrics}

For default prediction, accuracy, precision, recall, and F1 are reported at
threshold 0.5. ROC-AUC and precision-recall AUC evaluate ranking without
requiring this threshold. For fairness, demographic parity difference and
equal opportunity difference compare approval rates and approval rates among
non-defaulters. The signed differences retain direction, while absolute
differences describe disparity size. Bootstrap intervals quantify uncertainty
and paired intervals quantify changes between model conditions.

\section{Baselines and Comparison Systems}

The pre-thesis comparison provides a reference under its original protocol:
logistic regression achieved ROC-AUC 0.8632, random forest 0.9207, XGBoost
0.9388, and LightGBM 0.9431. The LightGBM margin over XGBoost is only 0.004
on that single split and is not treated as statistically significant. The
principal comparisons in this thesis are instead M0 versus M1 and each
unweighted model versus its reweighed version. This isolates the effects of
age removal and reweighing.
